\documentclass{article}

\usepackage{amsmath}
\usepackage{amssymb}
\usepackage{hyperref}

\title{SNIFs: Safe Native Implemented Functions for the BEAM via WebAssembly Sandboxing}
\author{Jonathan D. A. Jewell}
\date{2026}

\begin{document}

\maketitle

\begin{abstract}
Native Implemented Functions (NIFs) in the BEAM virtual machine enable integration with native code, but any fault within a NIF terminates the entire VM. This paper examines a practical approach to enforcing crash isolation by compiling native code to WebAssembly and executing it within a sandboxed runtime.

Using Zig targeting WASM and the \texttt{wasmex} interface, we show that a range of failure modes—including out-of-bounds access, explicit panics, unreachable instructions, integer overflow, and divide-by-zero—are trapped and returned as \texttt{\{:error, reason\}} tuples without terminating the BEAM VM.

We formalise a crash-isolation property under the operational semantics of WebAssembly and validate it empirically using an integration test suite (11/11 tests passing). In addition, the accompanying artifact includes machine-checked verification of core interface invariants (ABI correctness and API type safety) using Idris2 and Lean~4 (7 proofs total).

These results indicate that WebAssembly sandboxing provides a viable mechanism for isolating native faults in BEAM systems.
\end{abstract}

\section{Introduction}

The BEAM virtual machine provides a robust concurrency model, but Native Implemented Functions (NIFs) bypass many of its safety guarantees. A crash within a NIF typically terminates the entire VM, making NIFs a significant reliability risk.

This work evaluates an approach in which native code is compiled to WebAssembly and executed within a sandboxed runtime. Instead of propagating faults to the host, WebAssembly traps are captured and translated into error values at the BEAM level.

The contribution of this paper is not a new virtual machine or language extension, but a concrete demonstration that existing WebAssembly tooling can be used to enforce crash isolation for NIF-style integrations, together with a reproducible artifact and supporting formalisation.

\section{Architecture}

The SNIF architecture consists of three components:

\begin{itemize}
  \item Native functions written in Zig and compiled to WebAssembly.
  \item A WebAssembly runtime (Wasmtime via \texttt{wasmex}) embedded within the BEAM.
  \item A translation layer mapping traps to BEAM-level error values.
\end{itemize}

All guest execution occurs within WebAssembly linear memory. Faults such as out-of-bounds access or explicit traps are handled by the runtime and do not directly affect the host process.

A critical requirement is the use of safe compilation settings. In particular, Zig must be compiled with \texttt{-OReleaseSafe}; unsafe modes remove bounds checks and invalidate the safety guarantees relied upon here.

\section{Evaluation}

We evaluate the approach using a set of representative failure modes:

\begin{itemize}
  \item Out-of-bounds memory access
  \item Explicit panic
  \item Unreachable instruction
  \item Integer overflow
  \item Divide-by-zero
\end{itemize}

All failure cases are exercised via integration tests within a BEAM application. In all cases, the WebAssembly runtime traps and returns an error value without terminating the BEAM VM.

Across the test suite, 11/11 integration tests pass under safe compilation settings.

\section{Crash Isolation Property}

We describe the following property informally:

\textbf{Crash Isolation.} For any SNIF call executed via the WebAssembly runtime, any fault in guest execution results in a trapped execution state that is translated into a BEAM-level error value, without terminating the host VM.

This property relies on the trapping semantics of WebAssembly as implemented by the runtime (Wasmtime). Under these semantics, faults in guest execution do not escape the sandboxed environment.

\section{Formal Verification}

The accompanying artifact includes machine-checked verification of key interface invariants using Idris2 and Lean~4.

These proofs cover:

\begin{itemize}
  \item ABI correctness (data layout and calling conventions)
  \item Memory safety invariants (pointer validity and bounds constraints)
  \item API type safety (correctness of SNIF result representations)
\end{itemize}

All core proofs (7 in total) are complete and mechanically verified. These artifacts provide an executable specification of the SNIF interface, complementing the informal reasoning presented in this paper.

\section{Limitations}

This work relies on the correctness of the WebAssembly runtime. Bugs in the runtime or violations of its safety guarantees would invalidate the isolation property.

Additionally, unsafe compilation modes remove necessary checks and can lead to silent incorrect results rather than traps.

Performance is not evaluated in detail here; the focus is on correctness and isolation rather than throughput.

\section{Conclusion}

This paper demonstrates that WebAssembly sandboxing can be used to enforce crash isolation for NIF-style integrations in the BEAM. By combining safe compilation, runtime trapping semantics, and a translation layer, native faults can be contained and represented as ordinary error values.

The accompanying artifact provides a complete, reproducible implementation together with formal verification of key invariants, supporting the use of this approach in practice.

\bibliographystyle{plain}
\bibliography{references}

\end{document}
